% TEMPLATE-MILC-v3.tex - plantilla maestra UNICA de las guias MILC v3.
%
% La usan las cuatro familias de guias, cada una con su driver:
%   · Grados 6-11          -> scripts/build-guias-g11.py
%   · Bebras               -> scripts/build-guias-bebras.py
%   · Semillero            -> scripts/build-guias-semillero.py
%   · Territorio interior  -> scripts/build-guias-territorio-interior.py
%
% Antes habia tres copias casi identicas (~400 lineas iguales, 6 distintas):
% cada mejora habia que aplicarla tres veces, y en la practica no se hacia
% (el motor de recursos vivia solo en dos de ellas). Lo que varia entre
% familias esta parametrizado en 6 placeholders que cada driver llena
% (se nombran aqui SIN su sintaxis real: el builder reemplaza por texto plano,
% tambien dentro de los comentarios, y un valor multilinea rompe el preambulo):
%
%   HEADER_IZQ        encabezado izquierdo de pagina
%   HEADER_DER        encabezado derecho de pagina
%   PORTADA_ETIQUETA  rotulo sobre el numero grande (GRADO/MOMENTO/MODULO)
%   PORTADA_SERIE     linea de serie de la portada
%   PORTADA_ID        identificador de la guia en la portada
%   PORTADA_CIRCULO   el circulito de grado (°), o vacio si no aplica
%   PDF_TITULO        titulo en texto plano (sin \textbf ni comillas TeX) para
%                     los metadatos del PDF (/Title); lo produce lib_xelatex.texto_pdf
%
% Composicion (auditoria 2026-09-03): las bandas de estacion piden espacio
% (\Needspace*) y prohiben el salto justo despues (\nopagebreak), asi no quedan
% huerfanas al pie; las cajas largas (softbox, drawbox, infoband, puentebox) son
% `breakable`, asi una caja de media pagina ya no arrastra todo a la siguiente
% dejando la anterior al 40 %. Todo texto sobre mostaza o verde va en milcNegro
% (blanco sobre #E5B400 da 1,93:1 y sobre #58A923 2,95:1; ninguno pasa WCAG AA).
% El resto de claves las documenta content/guias/_SCHEMA.md. El script valida
% que no queden placeholders sin reemplazar antes de escribir el archivo final.

% Etiquetado PDF (latex-lab): /Lang, estructura y texto alternativo de las
% figuras, para lectores de pantalla. Debe ir ANTES de \documentclass.
\DocumentMetadata{lang=es-CO,tagging=on}
\documentclass[11pt,letterpaper]{article}
% headheight=24pt: la cabecera derecha lleva dos lineas (identidad de la guia
% y estacion actual). headsep se acorta en la misma medida para que el bloque
% de texto quede exactamente donde estaba con headheight=12pt.
\usepackage[margin=2.0cm,headheight=24pt,headsep=13pt]{geometry}
\usepackage[table]{xcolor}
\usepackage{fontspec}
\usepackage[spanish]{babel}
\usepackage{tikz}
% Librerías TikZ para los diagramas de las guías (bloque `recursos:`):
%  · babel  → babel en español vuelve activos caracteres como «>», y TikZ los usa
%             en sus opciones (>=stealth, ->). Sin esta librería, cualquier
%             diagrama con flechas revienta con «\language@active@arg> has an
%             extra }». Debe ir ANTES de que se dibuje nada.
%  · positioning        → sintaxis «below left=2pt and 3pt».
%  · decorations.pathreplacing → llaves ({brace}) para señalar rangos.
\usetikzlibrary{babel,positioning,decorations.pathreplacing}
\usepackage{graphicx}
% Circuitos: el trabajo de electrónica se hace hoy con micro:bit y sus sensores
% integrados (MakeCode, sin cableado), así que no se necesitan esquemáticos.
% Si algún día entran montajes con protoboard, basta descomentar esta línea:
% los diagramas .tex/.tikz declarados en `recursos:` ya se incrustan solos.
% \usepackage{circuitikz}
\usepackage[most]{tcolorbox}
\usepackage{tabularx}
\usepackage{array}
\usepackage{fancyhdr}
\usepackage{titlesec}
\usepackage[document]{ragged2e}  % bandera a la derecha en todo el cuerpo
\usepackage{enumitem}
\usepackage{microtype}
\usepackage{etoolbox}
\usepackage{needspace}
% hyperref al final del preambulo: metadatos del PDF (titulo, autor, idioma,
% familia), marcadores por estacion (\pdfbookmark) y enlaces sin recuadro.
\usepackage[hidelinks,
  pdftitle={Revolución industrial — máquinas que cambiaron el trabajo},
  pdfauthor={PhD. Álvaro Cárdenas Orozco},
  pdfsubject={Tecnología e Informática · Grado 9 · Guía 5 · MILC v3},
  pdflang={es-CO},
  bookmarksopen=true,bookmarksnumbered=false]{hyperref}

% Helvetica Neue no trae la flecha U+2192, asi que cada «→» escrito literalmente
% en el YAML se compilaba a NADA, en silencio (462 flechas en 61 guias: rutas del
% tipo «paleta → tipografia → lockup» salian rotas). Mapeamos el caracter a su
% version matematica, que si tiene glifo. Asi el autor puede seguir escribiendo
% «→» comodamente en el YAML.
\usepackage{newunicodechar}
\newunicodechar{→}{\ensuremath{\rightarrow}}
% Mismo problema con los simbolos de marca y vinetas: el «✓» de «una palomita ✓»
% salia como cuadrito vacio en el PDF de todas las guias que lo usaban.
\usepackage{pifont}
\newunicodechar{✓}{\ding{51}}
\newunicodechar{✗}{\ding{55}}
\newunicodechar{✔}{\ding{52}}
\newunicodechar{•}{\textbullet}
\newunicodechar{≥}{\ensuremath{\geq}}
\newunicodechar{≤}{\ensuremath{\leq}}

\setmainfont{Helvetica Neue}
\setsansfont{Helvetica Neue}
\setlength{\parindent}{0pt}
\setlength{\parskip}{6pt}
% Sin particion de palabras: en 6.º-7.º una palabra cortada por guion al final
% de linea («Comuni-cacion») frena la lectura mucho mas que una linea corta.
\hyphenpenalty=10000
\exhyphenpenalty=10000
\pretolerance=2000
\tolerance=3000
\setlength{\emergencystretch}{3em}
\linespread{1.10}
\renewcommand{\arraystretch}{1.25}
\setlist{leftmargin=5mm,itemsep=1mm,topsep=1mm}
\newcolumntype{Y}{>{\RaggedRight\arraybackslash}X}

\definecolor{milcMagenta}{HTML}{E6007E}
\definecolor{milcVino}{HTML}{5A0038}
\definecolor{milcTurquesa}{HTML}{0093A5}
\definecolor{milcVerde}{HTML}{58A923}
\definecolor{milcMostaza}{HTML}{E5B400}
\definecolor{milcOcre}{HTML}{B86600}
\definecolor{milcNegro}{HTML}{241816}
\definecolor{milcGris}{HTML}{F7F7F4}
\definecolor{milcLinea}{HTML}{E8E2DD}
\definecolor{milcGradePrimary}{HTML}{7C3AED}
\definecolor{milcGradeDark}{HTML}{4C1D95}
\definecolor{milcGradeSoft}{HTML}{EDE9FE}
\definecolor{milcGradeAccent}{HTML}{CFE7FF}
\definecolor{milcGradeText}{HTML}{FFFFFF}
\color{milcNegro}

\pagestyle{fancy}
\fancyhf{}
% Cabecera de dos lineas: identidad de la guia arriba y, a la derecha y en
% gris, la estacion en curso (\markright desde \milcestacion). Antes de la
% primera estacion la segunda linea va vacia; el \strut mantiene la altura
% para que la linea de arriba no baile entre paginas.
\lhead{\small Tecnología e Informática · Grado 9\\[-1pt]{\footnotesize\strut}}
\rhead{\small Guía 5 · MILC v3\\[-1pt]{\footnotesize\strut\color{milcNegro!65}\rightmark}}
\cfoot{\small\thepage}
\renewcommand{\headrulewidth}{0pt}

% Sobre mostaza (#E5B400) y verde (#58A923) el texto va en milcNegro; sobre el
% resto de la paleta, en blanco. Se usa dentro de fonttitle/fontupper, donde un
% \color posterior manda sobre coltitle.
\newcommand{\milctintasobre}[1]{%
  \ifboolexpr{test{\ifstrequal{#1}{milcMostaza}} or test{\ifstrequal{#1}{milcVerde}}}%
    {\color{milcNegro}}{\color{white}}}

\newtcolorbox{titlebox}[1]{
  enhanced,colback=#1,colframe=#1,arc=7mm,boxrule=0pt,
  left=7mm,right=7mm,top=3mm,bottom=3mm,
  before skip=8pt,after skip=8pt,
  fontupper=\sffamily\bfseries\Large\color{white}
}

\newtcolorbox{softbox}[3]{
  enhanced,breakable,pad at break=3mm,
  colback=#2,colframe=#1,borderline west={4pt}{0pt}{#1},
  arc=2mm,boxrule=.4pt,left=4mm,right=4mm,top=3mm,bottom=3mm,
  before skip=6pt,after skip=8pt,title={#3},coltitle=white,
  colbacktitle=#1,fonttitle=\sffamily\bfseries\milctintasobre{#1},fontupper=\RaggedRight,
  attach boxed title to top left={xshift=3mm,yshift=-2mm},
  boxed title style={arc=2mm, boxrule=0pt}
}

\newtcolorbox{drawbox}[2]{
  enhanced,breakable,pad at break=4mm,
  colback=white,colframe=#1,arc=4mm,boxrule=1pt,
  left=5mm,right=5mm,top=4mm,bottom=4mm,before skip=8pt,after skip=8pt,
  title={#2},coltitle=white,colbacktitle=#1,fonttitle=\sffamily\bfseries\milctintasobre{#1},
  fontupper=\RaggedRight,
  attach boxed title to top left={xshift=4mm,yshift=-2mm},
  boxed title style={arc=3mm, boxrule=0pt}
}

\newtcolorbox{infoband}[1]{
  enhanced,breakable,pad at break=2mm,colback=#1!8,colframe=#1!50,arc=2mm,boxrule=.5pt,
  left=4mm,right=4mm,top=2mm,bottom=2mm,before skip=4pt,after skip=4pt,
  fontupper=\small\RaggedRight
}

% Puente narrativo entre secciones (contrato editorial Conéctate).
% Da continuidad: "Acabas de X. Ahora vas a Y porque Z."
% Visualmente: banda izquierda en color del grado, texto en itálica pequeña.
\newtcolorbox{puentebox}{
  enhanced,breakable,pad at break=2mm,colback=milcGradeSoft,colframe=milcGradeDark,
  borderline west={3pt}{0pt}{milcGradeDark},
  arc=0mm,boxrule=0pt,left=5mm,right=4mm,top=2.5mm,bottom=2.5mm,
  before skip=4pt,after skip=8pt,
  fontupper=\itshape\small\color{milcNegro}\RaggedRight
}
% La flecha va en modo matematico a proposito: Helvetica Neue NO tiene el glifo
% U+2192, asi que \textrightarrow se compilaba a NADA («Missing character: There
% is no → in font Helvetica Neue») y el puente salia sin flecha. En modo
% matematico la flecha viene de la fuente matematica y si se dibuja.
\newcommand{\puente}[1]{%
  \begin{puentebox}%
    \textcolor{milcGradeDark}{$\rightarrow$}\hspace{2mm}#1%
  \end{puentebox}%
}

\newcommand{\linead}[1]{\noindent\color{milcLinea}\rule{#1}{0.5pt}\color{milcNegro}}
\newcommand{\respuesta}[1]{\par\vspace{2mm}\foreach \n in {1,...,#1}{\noindent\color{milcLinea}\rule{\linewidth}{0.5pt}\par\vspace{5mm}}\color{milcNegro}}
\newcommand{\checkbox}{\textcolor{milcGradeDark}{\fbox{\phantom{x}}}\hspace{5pt}}

% ─── Listas aireadas para las guías socioemocionales ────────────────────
% Componer las verificaciones como «\checkbox texto\\» deja el texto que salta
% de línea debajo del cuadrito y apelmaza el bloque. Estos entornos alinean la
% continuación y airean los ítems. Los usa Territorio Interior; cualquier
% familia puede hacerlo.
\newlist{checklista}{itemize}{1}
\setlist[checklista]{
  label=\textcolor{milcGradeDark}{\fbox{\phantom{x}}},
  leftmargin=9mm, labelsep=3.2mm, itemsep=2.4mm,
  topsep=2.5mm, parsep=0pt, partopsep=0pt, before=\RaggedRight
}
\newlist{pasoslista}{enumerate}{1}
\setlist[pasoslista]{
  label=\textcolor{milcGradeDark}{\sffamily\bfseries\arabic*.},
  leftmargin=8mm, labelsep=3mm, itemsep=2.8mm,
  topsep=1mm, parsep=0pt, partopsep=0pt, before=\RaggedRight
}

\newcommand{\phasecard}[4]{
\begin{tcolorbox}[enhanced,colback=#3,colframe=#2,arc=3mm,boxrule=.5pt,height=3.05cm,valign=center,halign=center,left=1.5mm,right=1.5mm,top=2mm,bottom=2mm]
{\sffamily\bfseries\fontsize{22}{24}\selectfont\textcolor{#2}{#1}}\par
{\sffamily\bfseries\small #4}
\end{tcolorbox}
}

% ── Piezas del rediseno 2026 (legibilidad 6.º-11.º) ───────────────────
% Banda de estacion: reemplaza el titlebox macizo. Titulo a la izquierda,
% dato practico (tiempo/modalidad) a la derecha, en una sola linea.
% Nunca huerfana: pide 9 lineas libres (o salta de pagina rellenando con
% \vfill) y prohibe el salto justo despues de la banda. Nueve y no seis:
% la caja partible que sigue necesita ~80pt (titulo adosado, relleno y dos
% lineas) para arrancar en la misma pagina; con menos, tcolorbox la manda
% entera a la siguiente y la banda se queda sola. El penalty 10000 va
% en `after`, ANTES del espacio posterior: un `after skip` (aunque sea 0pt)
% es un \vskip tras la caja, y ese glue es un punto de corte legal que un
% \nopagebreak escrito despues de \end{tcolorbox} ya no cierra. Cada banda
% deja marcador PDF y actualiza la cabecera derecha.
\newcounter{milcestacion}
\newcommand{\milcestacion}[2]{%
\Needspace*{9\baselineskip}%
\stepcounter{milcestacion}%
\pdfbookmark[1]{#2}{milcestacion\themilcestacion}%
\markright{#2}%
\begin{tcolorbox}[enhanced,colback=#1,colframe=#1,arc=2mm,boxrule=0pt,
  left=5mm,right=5mm,top=2.6mm,bottom=2.6mm,before skip=11pt,
  after={\par\nopagebreak\vskip7pt}]
{\sffamily\bfseries\fontsize{15}{18}\selectfont\milctintasobre{#1}#2}
\end{tcolorbox}}

% Tarjeta de paso numerada: sustituye las tablas «Pilar / Aplicacion», que
% eran formato de consulta docente, no de estudiante. El numero va en un
% circulo sobre el borde para que la secuencia se lea de un vistazo.
\newcommand{\milcpaso}[3]{%
\Needspace{4\baselineskip}%
\begin{tcolorbox}[enhanced,colback=white,colframe=milcLinea,arc=1.5mm,boxrule=.9pt,
  left=14mm,right=4mm,top=3mm,bottom=3mm,before skip=4pt,after skip=4pt,
  fontupper=\RaggedRight,
  overlay={\node[anchor=north west,circle,fill=#2,text=white,inner sep=0pt,
    minimum size=7.4mm,font=\sffamily\bfseries\small]
    at ([xshift=3.6mm,yshift=-3.2mm]frame.north west) {#1};}]
#3
\end{tcolorbox}}

% Casilla marcable de tamano util con lapiz (la anterior era \fbox{\phantom{x}},
% de alto variable segun el texto que la seguia).
\newcommand{\milcmarca}{\framebox[3.8mm]{\rule{0pt}{3.4mm}}}

% Fila marcable de lista suelta (checks de la estacion 1).
\newcommand{\milccheckrow}[1]{%
\par\vspace{1.8mm}\noindent
\makebox[6.4mm][l]{\raisebox{-.55mm}{\textcolor{milcNegro}{\milcmarca}}}%
\parbox[t]{\dimexpr\linewidth-6.4mm\relax}{\RaggedRight #1}\par}

% Celda marcable dentro de la rubrica (tabla de autoevaluacion). Misma
% sangria francesa, pero pensada para vivir dentro de una columna X.
\newcommand{\milcopcion}[1]{%
\makebox[5.2mm][l]{\raisebox{-.55mm}{\textcolor{milcNegro}{\milcmarca}}}%
\parbox[t]{\dimexpr\linewidth-5.2mm\relax}{\RaggedRight #1}}

% ── Recursos visuales de la guía ──────────────────────────────────────
% Declarados en el bloque `recursos:` del YAML y emitidos por el builder.
% \guiaFigura   → imagen rasterizada o PDF (foto, carta celeste, montaje).
% \guiaDiagrama → fuente TikZ/circuitikz (.tex) incrustada como vector:
%                 es la vía para circuitos y esquemáticos editables.
% [#1] = texto alternativo (alt) · #2 = ruta relativa al .tex · #3 = pie
% El alt llega al PDF etiquetado (/Alt) y lo lee el lector de pantalla.
\newcommand{\guiaFigura}[3][]{%
\begin{center}
\includegraphics[width=0.88\linewidth,height=0.38\textheight,keepaspectratio,alt={#1}]{#2}\\[1.5mm]
{\footnotesize\itshape\textcolor{milcNegro!75}{#3}}
\end{center}
}

\newcommand{\guiaDiagrama}[2]{%
\begin{center}
\input{#1}\\[1.5mm]
{\footnotesize\itshape\textcolor{milcNegro!75}{#2}}
\end{center}
}

\begin{document}
\thispagestyle{empty}

% ── Portada ───────────────────────────────────────────────────────────
% Antes: un rectangulo de color a pagina completa. Se veia bien en pantalla,
% pero en la fotocopiadora del colegio se comia el toner y salia gris sucio.
% Ahora la banda ocupa el tercio superior y el resto va en blanco.
\begin{tikzpicture}[remember picture,overlay]
  \path[fill=milcGradePrimary]
    (current page.north west) rectangle ([yshift=-10.6cm]current page.north east);

  \node[anchor=north west,text=milcGradeText] at ([xshift=2.0cm,yshift=-1.55cm]current page.north west)
    {\sffamily\bfseries\fontsize{12}{14}\selectfont GRADO};
  \node[anchor=north west,text=milcGradeText,opacity=.85] at ([xshift=2.0cm,yshift=-2.35cm]current page.north west)
    {\sffamily\bfseries\fontsize{8.5}{10}\selectfont SERIE GUÍAS · TECNOLOGÍA E INFORMÁTICA};
  \node[anchor=north east,text=milcGradeText,opacity=.85,align=right] at ([xshift=-2.0cm,yshift=-1.55cm]current page.north east)
    {\sffamily\bfseries\fontsize{9}{12}\selectfont I.E. Sor María Juliana\\Cartago, Valle del Cauca};

  \node[anchor=west,text=milcGradeText] (gradonum) at ([xshift=1.85cm,yshift=-7.1cm]current page.north west)
    {\sffamily\bfseries\fontsize{112}{112}\selectfont 9};
  % El circulito de grado (°) solo aplica a las guias de grado («11°»); en
  % Bebras y Semillero el numero grande es un momento/modulo, no un grado.
  % Se posiciona relativo al nodo `gradonum`, no en coordenadas absolutas.
  \draw[milcGradeText,line width=1.6pt]
    ([xshift=2.0mm,yshift=-4.5mm]gradonum.north east) circle (.15cm);

  \node[anchor=west,text=milcGradeText,text width=11.4cm,align=left] at ([xshift=1.5cm,yshift=0.5cm]gradonum.east)
    {
     {\sffamily\bfseries\fontsize{9}{11}\selectfont\MakeUppercase{Período 1 · Historia de la técnica}}\\[3mm]
     {\sffamily\bfseries\fontsize{21}{25}\selectfont\setlength{\baselineskip}{27pt}Revolución industrial ---\\[4pt]máquinas que\\[4pt]cambiaron el trabajo}\\[3.5mm]
     {\sffamily\bfseries\fontsize{15}{18}\selectfont Guía 5-9-TIC}
    };
\end{tikzpicture}

\vspace*{9.4cm}
\vfill

{\sffamily\bfseries\fontsize{8}{10}\selectfont\textcolor{milcGradeDark}{LO QUE TE LLEVAS HOY}}

\begin{tcolorbox}[enhanced,colback=white,colframe=milcNegro,arc=2mm,boxrule=1.6pt,
  left=5mm,right=5mm,top=4mm,bottom=4mm,before skip=3pt,after skip=12pt,
  fontupper=\RaggedRight\fontsize{11}{15.5}\selectfont]
Un mapa visual de oficios de tu región: tres que la máquina sustituyó y tres que la máquina creó, cada uno con nombre concreto, fecha aproximada y lugar. Al menos uno sale de lo que te contó una persona mayor de tu familia. El mapa se cierra con un cuadro que nombra a quién le sirvió cada cambio y a quién no.
\end{tcolorbox}

\begin{tcolorbox}[enhanced,colback=milcGris,colframe=milcGris,arc=2mm,boxrule=0pt,
  left=5mm,right=5mm,top=3.5mm,bottom=3.5mm,after skip=12pt]
{\sffamily\bfseries\fontsize{8}{10}\selectfont\textcolor{milcNegro!55}{ESTA GUÍA ES DE}}\\[3mm]
\begin{tabularx}{\linewidth}{X p{3.2cm} p{3.2cm}}
\linead{\linewidth} & \linead{3.0cm} & \linead{3.0cm}\\[-1mm]
{\fontsize{8.5}{10}\selectfont\textcolor{milcNegro!55}{Nombre completo}} &
{\fontsize{8.5}{10}\selectfont\textcolor{milcNegro!55}{Grupo}} &
{\fontsize{8.5}{10}\selectfont\textcolor{milcNegro!55}{Fecha}}\\
\end{tabularx}
\end{tcolorbox}

\vfill

\noindent\textcolor{milcLinea}{\rule{\linewidth}{1pt}}\\[2mm]
{\sffamily\bfseries\fontsize{9.5}{12}\selectfont Autor: Dr. Álvaro Cárdenas Orozco}\\[1mm]
{\sffamily\fontsize{8.5}{11}\selectfont\textcolor{milcNegro!55}{Metodología MILC · apertura ancestral · pensamiento computacional · triángulo Dussel–estoicismo–Floridi}}

\newpage

\section*{Apertura: Revolución industrial --- máquinas que cambiaron el trabajo}

\begin{softbox}{milcGradeDark}{milcGradeSoft}{Saber ancestral}
En la finca cafetera del siglo XX los oficios venían asignados. Un informe consular británico de comienzos de siglo dice que la recolección «la llevan a cabo mujeres y niños». Ellas eran además escogedoras en la trilladora, cocineras de los trabajadores, encargadas de los animales, de la huerta y del pancoger (Ramírez Bacca, 2015). A la recolectora la llamaron \textbf{chapolera}, por la chapola, el arbolito de café recién nacido. Cien años después la desigualdad no es pasado. El DANE (2022) muestra que en los predios rurales de un solo dueño el 63,7~\% son hombres y el 36,3~\% mujeres. En ningún departamento del país hay paridad. Y hay un dato que incomoda más: hay más mujeres con título que mujeres tomando las decisiones productivas. La \textbf{cara de exclusión} es el ancla misma. Hoy la chapolera circula como disfraz folclórico para turistas, y ese disfraz tapa que el reparto de oficios fue una decisión, no una costumbre.\par\smallskip{\footnotesize\textcolor{milcNegro!70}{\textbf{Fuente.} Ramírez Bacca, R. (2015). Mujeres en la caficultura tradicional colombiana, 1910--1970. \emph{Historia y Memoria}, (10), 43--73. https://doi.org/10.19053/20275137.3200}}
\end{softbox}

\begin{softbox}{milcTurquesa}{milcGris}{Saber contemporáneo}
Cuando entra una máquina, los oficios de un lugar quedan en cuatro estados. \textbf{Sustituido}: la máquina hace lo que hacía una persona y el oficio desaparece. \textbf{Transformado}: el oficio sigue existiendo pero cambia de contenido. \textbf{Creado}: aparece un oficio que antes no tenía sentido. \textbf{Intacto}: la máquina no lo tocó, casi siempre porque no le salía a cuenta. Leer un cambio técnico es poner cada oficio en uno de esos cuatro estados y preguntar por el reparto. Quién ganó tiempo, quién ganó dinero y quién perdió el trabajo casi nunca son la misma persona. La revolución industrial no fue solo carbón y vapor: fue un reparto nuevo de quién hace qué. Y la parte que menos se cuenta es que ese reparto no cayó del cielo. Alguien decidió a quién le tocaba cada cosa.
\end{softbox}

\begin{softbox}{milcMostaza}{milcGris}{Pregunta-puente}
\textit{En la finca cafetera los oficios venían asignados y nadie los había escogido. ¿Qué oficio de tu casa está asignado hoy sin que nadie lo haya decidido en voz alta?}
\end{softbox}

\begin{softbox}{milcVerde}{milcGris}{Saber y saber hacer}
Al terminar podrás: (1) \textbf{identificar} oficios que desaparecieron y aparecieron en tu región, preguntándole a alguien que los vivió; (2) \textbf{explicar} los cuatro estados en que queda un oficio cuando entra una máquina; (3) \textbf{crear} un mapa de oficios ganados y perdidos que diga a quién le sirvió cada cambio.
\end{softbox}



\small
\begin{tabularx}{\textwidth}{>{\bfseries}p{3.0cm}Y}
\rowcolor{milcGradePrimary!12}Autor & Dr. Álvaro Cárdenas Orozco\\
DBA & Análisis crítico de la historia de la tecnología y su impacto social (MEN, T\&I)\\
Referentes & Dussel (técnica situada) · Estoicismo (téchne del cuerpo) · Floridi (infoesfera)\\
Producto final & Un mapa visual de oficios de tu región: tres que la máquina sustituyó y tres que la máquina creó, cada uno con nombre concreto, fecha aproximada y lugar. Al menos uno sale de lo que te contó una persona mayor de tu familia. El mapa se cierra con un cuadro que nombra a quién le sirvió cada cambio y a quién no.\\
\end{tabularx}
\normalsize

\puente{Vienes de definir qué es técnica y de estudiar la balanza y las máquinas simples. Hoy miras qué pasa cuando la máquina entra a un lugar donde ya había gente trabajando. En la sesión 6 llega la imprenta, que hizo lo mismo con la palabra escrita.}

\milcestacion{milcGradeDark}{Ruta MILC: así vamos a aprender}
\begin{tabularx}{\textwidth}{>{\centering\arraybackslash}X>{\centering\arraybackslash}X>{\centering\arraybackslash}X>{\centering\arraybackslash}X}
\phasecard{1}{milcVerde}{milcGris}{Escucha\\[-1mm]\small Observo, escucho y pregunto.} &
\phasecard{2}{milcTurquesa}{milcGris}{Entender\\[-1mm]\small Sistematizo y ordeno.} &
\phasecard{3}{milcMagenta}{milcGris}{Hacer\\[-1mm]\small Construyo y aplico.} &
\phasecard{4}{milcVino}{milcGris}{Cierre\\[-1mm]\small Evalúo y reflexiono.}
\end{tabularx}


\puente{Antes de la teoría vas a preguntar. Una llamada a alguien mayor de tu familia vale más que diez páginas de libro, porque esa persona vio desaparecer un oficio con nombre propio.}

\milcestacion{milcVerde}{Estación 1 de 3 · Escucha}

\begin{softbox}{milcVerde}{milcGris}{Escena para escuchar}
\textbf{Actividad 1 · IDENTIFICA --- Un oficio que ya no existe} (15 min · individual). (1) Llama o habla con alguien de tu familia mayor de cuarenta años. (2) Pregúntale qué oficio había en su juventud que hoy ya no existe. (3) Pregúntale qué oficio existe hoy que no existía entonces. (4) Anota el nombre exacto de cada oficio, el año aproximado y el lugar. (5) Pregúntale qué le pasó a la gente que vivía del oficio que desapareció.
\end{softbox}

{\sffamily\bfseries\fontsize{8}{10}\selectfont\textcolor{milcNegro!55}{MARCA CUANDO LO HAYAS HECHO}}
\milccheckrow{Tengo el nombre exacto de un oficio desaparecido, con año y lugar.}
\milccheckrow{Tengo el nombre de un oficio nuevo, con año aproximado.}
\milccheckrow{Anoté qué le pasó a la gente que vivía del oficio que desapareció.}

\begin{infoband}{milcVerde}


\textbf{Lo que va al cuaderno (Actividad 1 · IDENTIFICA).} \textbf{Título:} «Actividad 1 --- Un oficio que ya no existe». \textbf{Formato:} ficha con el nombre de quien te contó, el oficio desaparecido con año y lugar, el oficio nuevo, y qué pasó con la gente. \textbf{Extensión:} un tercio de página. \textbf{Sabes que terminaste cuando} los dos oficios tienen nombre propio, no categoría.
\end{infoband}

\puente{Ya tienes testimonio. Ahora vas a ordenarlo en los cuatro estados y a mirar el reparto: quién ganó y quién perdió con cada cambio.}

\milcestacion{milcTurquesa}{Estación 2 de 3 · Entender}

\begin{softbox}{milcTurquesa}{milcGris}{Lo que vas a entender}
\textbf{Actividad 2 · EXPLICA --- Los cuatro estados de un oficio} (20 min · parejas). (1) Con tu pareja, escriban los cuatro estados con una frase propia cada uno. (2) Clasifiquen en ellos los oficios que trajeron de sus entrevistas. (3) Elijan uno y escriban quién ganó tiempo, quién ganó dinero y quién perdió el trabajo. (4) Comprueben si las tres respuestas son la misma persona y escriban qué significa que no lo sean.
\end{softbox}

\milcpaso{1}{milcTurquesa}{\textbf{Los cuatro estados.} Sustituido: la máquina hace lo que hacía una persona. Transformado: el oficio sigue, pero cambia por dentro. Creado: aparece uno que antes no tenía sentido. Intacto: la máquina no lo tocó, casi siempre porque no le salía a cuenta.}
\milcpaso{2}{milcTurquesa}{\textbf{El reparto es la pregunta.} Quién gana tiempo, quién gana dinero y quién pierde el trabajo casi nunca son la misma persona. Cuando lo son, no hubo conflicto. Cuando no lo son, hubo un reparto y conviene nombrarlo.}
\milcpaso{3}{milcTurquesa}{\textbf{Los oficios asignados no son costumbre.} En la finca cafetera alguien decidió que la recolección y la escogida eran de mujeres y niños. No fue la naturaleza ni la tradición: fue una decisión que se repitió tanto que dejó de parecer una.}
\milcpaso{4}{milcTurquesa}{\textbf{Tener el título no es decidir.} El dato del DANE lo muestra: hay más mujeres con predio a su nombre que mujeres tomando las decisiones productivas. Poseer y decidir son dos cosas distintas, y confundirlas esconde la desigualdad.}

\begin{drawbox}{milcTurquesa}{Un oficio cuando entra una máquina}
\textbf{Sustituido:} la máquina lo hace. \\[1mm] \textbf{Transformado:} sigue, pero cambia por dentro. \\[1mm] \textbf{Creado:} antes no tenía sentido. \\[1mm] \textbf{Intacto:} no salía a cuenta automatizarlo. \\[1mm] \textbf{Pregunta:} ¿quién ganó tiempo, quién ganó dinero, quién perdió el trabajo?
\end{drawbox}

\begin{softbox}{milcMostaza}{milcGris}{Errores comunes y su corrección}
(1) \textbf{Decir «el progreso» como responsable}: el progreso no decide, deciden personas e instituciones. (2) \textbf{Nombrar categorías en vez de oficios}: «los obreros» no es un oficio; «trilladora» o «telegrafista» sí. (3) \textbf{Tratar la desigualdad como pasado}: la brecha de propiedad rural es un dato de 2022. (4) \textbf{Romantizar el oficio perdido}: muchos eran durísimos y mal pagos. (5) \textbf{Olvidar los oficios intactos}: dicen mucho sobre qué no le convino automatizar a nadie.
\end{softbox}

\begin{infoband}{milcTurquesa}


\textbf{Lo que va al cuaderno (Actividad 2 · EXPLICA).} \textbf{Título:} «Actividad 2 --- Los cuatro estados de un oficio». \textbf{Formato:} los cuatro estados con frase propia, los oficios de las entrevistas clasificados y el reparto de uno de ellos en tres renglones. \textbf{Extensión:} media página. \textbf{Sabes que terminaste cuando} sabes decir si las tres respuestas del reparto son la misma persona.
\end{infoband}

\puente{Tienes el criterio. Toca dibujar el mapa de tu región y sostener con datos quién salió ganando.}

\milcestacion{milcMagenta}{Estación 3 de 3 · Hacer}

\begin{softbox}{milcMagenta}{milcGris}{Cómo vas a construir tu producto}
\textbf{Actividad 3 · CREA --- Mapa de oficios ganados y perdidos} (30 min · individual). (1) Dibuja una línea de tiempo de tu región, del siglo XX a hoy. (2) Pon tres oficios que la máquina sustituyó, con nombre, año aproximado y lugar. (3) Pon tres oficios que la máquina creó, con los mismos datos. (4) Marca cuál de los seis salió de tu entrevista. (5) Cierra con un cuadro que diga, para dos de ellos, a quién le sirvió el cambio y a quién no. (6) Muéstraselo a un compañero sin explicarle nada y anota qué no se entendió.
\end{softbox}

\milcpaso{1}{milcMagenta}{\textbf{Tu entrega tiene tres piezas.} La línea de tiempo con seis oficios. La marca en el que salió de tu entrevista. El cuadro del reparto para dos de ellos.}
\milcpaso{2}{milcMagenta}{\textbf{Nombre propio, año y lugar.} «Telegrafista, años sesenta, oficina de Cartago» es un dato. «Los trabajos de antes» no lo es. Sin nombre propio no se puede comprobar nada.}
\milcpaso{3}{milcMagenta}{\textbf{El mapa se lee sin ti al lado.} Si tu compañero necesita que se lo expliques, todavía falta jerarquía visual o falta un rótulo. Eso se arregla antes de entregar, no después.}
\milcpaso{4}{milcMagenta}{\textbf{El cuadro del reparto es lo difícil.} Escribir «a quién le sirvió» obliga a nombrar a alguien: una empresa, unos dueños, unos clientes. Si no logras nombrar a nadie, es que todavía no averiguaste lo suficiente.}

\begin{drawbox}{milcMagenta}{Antes de entregar}
\checkbox Tres oficios sustituidos, con nombre, año y lugar. \\[1mm] \checkbox Tres oficios creados, con los mismos datos. \\[1mm] \checkbox Marcado el que salió de tu entrevista. \\[1mm] \checkbox Cuadro del reparto para dos oficios. \\[1mm] \checkbox El cuadro nombra a alguien concreto, no «el progreso». \\[1mm] \checkbox Un compañero lo entendió sin explicación.
\end{drawbox}

\begin{softbox}{milcMagenta}{milcGris}{Plantilla de trabajo}
\textbf{Oficio perdido.} \textit{Nombre, año, lugar.} \\[1mm] \textbf{Oficio creado.} \textit{Nombre, año, lugar.} \\[1mm] \textbf{De la entrevista.} \textit{Quién me lo contó.} \\[1mm] \textbf{Reparto.} \textit{Quién ganó tiempo, quién ganó dinero, quién perdió el trabajo.} \\[1mm] \textbf{Prueba.} \textit{Qué no entendió mi compañero.}
\end{softbox}

\begin{infoband}{milcMagenta}


\textbf{Lo que va al cuaderno (Actividad 3 · CREA).} \textbf{Título:} «Actividad 3 --- Mapa de oficios ganados y perdidos». \textbf{Formato:} la línea de tiempo con los seis oficios rotulados y el cuadro del reparto de dos de ellos. \textbf{Extensión:} media página. \textbf{Sabes que terminaste cuando} un compañero leyó el mapa sin que le explicaras nada.
\end{infoband}

\puente{Lo que entregas hoy: el mapa con seis oficios y el cuadro del reparto. La prueba de calidad: alguien que no estuvo en clase lo entiende sin que se lo expliques.}

\milcestacion{milcMostaza}{Producto de la sesión}
\textbf{Mapa de seis oficios ganados y perdidos + cuadro del reparto}.
\begin{softbox}{milcMostaza}{milcGris}{Criterios de calidad}
(1) Tres oficios sustituidos y tres creados, todos con nombre propio. (2) Cada oficio tiene año aproximado y lugar concreto. (3) Al menos uno salió de la entrevista a una persona mayor de tu familia. (4) El cuadro del reparto nombra a alguien concreto, no a «el progreso». (5) Un compañero entendió el mapa sin explicación oral.
\end{softbox}

\puente{Antes de cerrar, mira lo que hiciste desde cinco lados. Un cambio técnico siempre reparte, y decir que «el progreso» es el responsable es no nombrar a nadie.}

\milcestacion{milcVino}{Evaluación liberadora · cinco dimensiones}

\begin{softbox}{milcVino}{milcGris}{Sentido de la evaluación}
Evaluar no es solo revisar una nota. Es reconocer qué aprendí, qué cambió en mí, qué emociones manejé, cómo me relaciono con la ciudad y la comunidad, y qué le dejaré a quien viene detrás.
\end{softbox}

\textbf{Mi reflexión liberadora en cinco dimensiones.} Completa cada frase iniciada:

\small
\begin{tabularx}{\textwidth}{>{\bfseries\RaggedRight\arraybackslash}p{3.4cm}Y>{\RaggedRight\arraybackslash}p{4.6cm}}
\rowcolor{milcVino}\color{white}Dimensión & \color{white}\bfseries Mi respuesta (frame starter) & \color{white}\bfseries Algo que descubrí\\
1. Personal & Lo que más cambió en mí fue \linead{6.5cm} & \linead{4.2cm}\\
2. Emocional & La emoción más fuerte fue \linead{2.5cm} y la regulé \linead{3cm} & \linead{4.2cm}\\
3. Ciudadana & En democracia esto importa porque \linead{6cm} & \linead{4.2cm}\\
4. Local (Cartago) & En mi barrio sería útil para \linead{6.5cm} & \linead{4.2cm}\\
5. Intergeneracional & A mi abuela/abuelo le diría \linead{6.5cm} & \linead{4.2cm}\\
\end{tabularx}
\normalsize

\begin{infoband}{milcVino}
\textbf{Para la próxima sustentación.} Vuelve a esta tabla en seis meses. Verás que las respuestas cambian — no porque lo aprendido cambie, sino porque tú cambias. La evaluación liberadora es una conversación contigo a través del tiempo.
\end{infoband}


\puente{Para terminar, tres citas verificadas sobre máquinas y trabajo: Dussel sobre la máquina que vuelve instrumento al hombre, Epicteto sobre las opiniones que atormentan, Floridi sobre una libertad que no ocurre en el vacío.}

\milcestacion{milcGradeDark}{Triángulo de pensamiento · cierre filosófico}

\begin{softbox}{milcVino}{milcGris}{Dussel · lente del nosotros}
\textit{«La máquina, el robot, el «rostro material» del capital ha hecho del «rostro del hombre» un instrumento de sí mismo.»}

{\footnotesize\itshape\color{milcNegro!70}--- Enrique Dussel · Filosofía de la liberación (1977), §2.5.6.3}

No es que la máquina sea mala. Es que, cuando el que decide es el que la compra, la persona pasa a medirse por lo que la máquina necesita de ella. La escogedora de la trilladora sabía eso mucho antes de que se escribiera.

\textbf{¿En qué parte de mi día me estoy adaptando yo a una máquina, en vez de al revés?}
\end{softbox}
\linead{\linewidth}\\[1mm]
\linead{\linewidth}

\begin{softbox}{milcMostaza}{milcGris}{Estoicismo · Epicteto · Enquiridión, 5 (c. 125 d.C.) · cuidado interior}
\textit{«No son las cosas las que atormentan a los hombres, sino las opiniones que se tienen de ellas.»}

{\footnotesize\itshape\color{milcNegro!70}--- Epicteto · Enquiridión, 5 (c. 125 d.C.)}

Que la recolección fuera «cosa de mujeres y niños» no venía del café ni de la máquina. Venía de una opinión repetida hasta parecer un hecho. Distinguir el hecho de la opinión heredada es la mitad del trabajo de hoy.

\textbf{¿Qué reparto de tareas doy por natural en mi casa sin que nadie lo haya decidido en voz alta?}
\end{softbox}
\linead{\linewidth}\\[1mm]
\linead{\linewidth}

\begin{softbox}{milcTurquesa}{milcGris}{Floridi · lente de la infoesfera}
\textit{«Nuestro yo es a la vez libre y social: la libertad no ocurre en el vacío, sino en un espacio de posibilidades y de límites. (trad. propia)»}

{\footnotesize\itshape\color{milcNegro!70}--- The Onlife Initiative (ed. Luciano Floridi) · The Onlife Manifesto (2015), § 4.2}

Nadie elige su oficio en el aire. Se elige entre lo que hay, con la máquina que ya llegó y el reparto que ya existe. Ver esos límites no quita libertad: es lo que permite usarla con puntería.

\textbf{¿Qué opciones de trabajo tengo de verdad hoy, y quién definió esa lista?}
\end{softbox}
\linead{\linewidth}\\[1mm]
\linead{\linewidth}

\begin{infoband}{milcNegro}
\textbf{Nota para el docente.} Las tres son citas textuales y llevan su referencia debajo. Puedes remitir a la obra en clase.
\end{infoband}

\milcestacion{milcVino}{Mi autoevaluación · marca una casilla por fila}

{\small
\setlength{\extrarowheight}{2.2pt}
\begin{tabularx}{\textwidth}{>{\RaggedRight\arraybackslash\bfseries}p{3.0cm}YYY}
\rowcolor{milcVino}\color{white}\bfseries Criterio & \color{white}\bfseries Lo logré & \color{white}\bfseries En proceso & \color{white}\bfseries Necesito apoyo\\[.6mm]
Comprensión del tema & \milcopcion{Explico las ideas con mis palabras.} & \milcopcion{Reconozco las ideas, falta precisión.} & \milcopcion{Necesito repasar lo básico.}\\[.8mm]
\rowcolor{milcGris}Pensamiento computacional & \milcopcion{Usé los pilares al armar mi producto.} & \milcopcion{Usé algunos pilares.} & \milcopcion{Necesito releer la estación 2.}\\[.8mm]
Aplicación y producto & \milcopcion{Mi producto cumple los criterios.} & \milcopcion{Está armado, con ajustes pendientes.} & \milcopcion{Aún está en construcción.}\\[.8mm]
\rowcolor{milcGris}Reflexión 5 dimensiones & \milcopcion{Respondí cada una con honestidad.} & \milcopcion{Respondí algunas con detalle.} & \milcopcion{Mis respuestas son muy generales.}\\[.8mm]
Triángulo de pensamiento & \milcopcion{Conecté las 3 voces con mi trabajo.} & \milcopcion{Conecté al menos una.} & \milcopcion{No las relaciono con mi proceso.}\\[.8mm]
\end{tabularx}}

\vspace{3mm}
\textbf{Hoy entendí que:}\\[1mm]
\linead{\linewidth}\\[1mm]
\linead{\linewidth}

\textbf{Mi compromiso para la próxima sesión es:}\\[1mm]
\linead{\linewidth}\\[1mm]
\linead{\linewidth}

\begin{softbox}{milcMostaza}{milcGris}{Cita de despedida · Marco Aurelio}
\textit{``El obstáculo en el camino se vuelve el camino. Lo que estorba al camino se vuelve camino.''}\\[1mm]
Cada error de hoy, bien observado, se convierte en pieza del camino que pisarás mañana.
\end{softbox}

\small
\begin{tabularx}{\textwidth}{>{\bfseries}p{3.6cm}Y}
\rowcolor{milcGradeDark!12}Nombre completo: & \linead{10cm}\\
Fecha: & \linead{4cm} \quad Curso/grupo: \linead{4cm}\\
Mi firma: & \linead{9cm}\\
\end{tabularx}
\normalsize

\begin{infoband}{milcGradeDark}
\textbf{Para la próxima sesión.} Llega con el mapa y el cuadro del reparto. Pregúntale a la misma persona que entrevistaste si el oficio nuevo paga mejor que el que desapareció. En la sesión 6 vas a trabajar la imprenta, que hizo con la palabra escrita lo que la máquina hizo con el trabajo.
\end{infoband}

\end{document}
